\documentclass[12pt, a4paper]{ctexart}
\usepackage{amsmath, amssymb, amsthm}
\usepackage{fancyhdr}
\usepackage{geometry}
\usepackage{enumitem}
\usepackage{graphicx}
\usepackage{tikz}
\usepackage{pgfplots}
\pgfplotsset{compat=1.18}

\geometry{left=2cm, right=2cm, top=2.5cm, bottom=2.5cm}

\setCJKmainfont{Hiragino Sans GB}
\setCJKsansfont{Heiti SC}

\pagestyle{fancy}
\fancyhf{}
\fancyhead[L]{高考数学满分突破·数学思想方法总结}
\fancyhead[R]{测试卷}
\fancyfoot[C]{\thepage}
\renewcommand{\headrulewidth}{0.4pt}

\setlength{\parindent}{0pt}
\setlength{\parskip}{6pt}

\newcommand{\score}[1]{\hfill\textbf{（#1分）}}

\begin{document}

\begin{center}
{\LARGE\bfseries 高考数学满分突破测试卷（十）}\\[8pt]
{\large 数学思想方法总结}\\[6pt]
考试时间：120分钟 \quad 满分：150分
\end{center}

\vspace{10pt}

{\large\bfseries 一、单项选择题：本题共8小题，每小题5分，共40分。每小题给出的四个选项中，只有一项是符合题目要求的。}

\vspace{8pt}

\begin{enumerate}

\item 设函数 $f(x)=x^3-3x+1$，则方程 $f(x)=0$ 在 $[-2,2]$ 上实根的个数为 \score{5}
\begin{enumerate}[label=\Alph*.]
    \item $1$
    \item $2$
    \item $3$
    \item $0$
\end{enumerate}

\item 用换元法，设 $t=\sqrt{x-1}$（$x\geq 1$），则方程 $x-3\sqrt{x-1}=1$ 化为 \score{5}
\begin{enumerate}[label=\Alph*.]
    \item $t^2-3t=0$
    \item $t^2-3t-2=0$
    \item $t^2+1-3t=1$，即 $t^2-3t=0$
    \item $t^2-3t+2=0$
\end{enumerate}

\item 设 $f(x)=\sin x+\cos x$，则 $f(x)$ 的最大值及取到最大值时 $x$ 的集合为 \score{5}
\begin{enumerate}[label=\Alph*.]
    \item 最大值为 $\sqrt{2}$，$x=\dfrac{\pi}{4}+2k\pi$，$k\in\mathbb{Z}$
    \item 最大值为 $2$，$x=\dfrac{\pi}{4}+2k\pi$，$k\in\mathbb{Z}$
    \item 最大值为 $\sqrt{2}$，$x=\dfrac{\pi}{2}+2k\pi$，$k\in\mathbb{Z}$
    \item 最大值为 $1$，$x=k\pi$，$k\in\mathbb{Z}$
\end{enumerate}

\item 设 $a_n=\dfrac{1}{n(n+1)}$，则 $\displaystyle\sum_{k=1}^{n}a_k$ 用裂项法求和的结果为 \score{5}
\begin{enumerate}[label=\Alph*.]
    \item $\dfrac{n}{n+1}$
    \item $\dfrac{n+1}{n+2}$
    \item $\dfrac{2n}{n+1}$
    \item $\dfrac{n}{2(n+1)}$
\end{enumerate}

\item 设 $f(x)=\ln x-x$，则 $f(x)$ 在 $(0,+\infty)$ 上的最大值为 \score{5}
\begin{enumerate}[label=\Alph*.]
    \item $0$
    \item $-1$
    \item $1$
    \item 不存在
\end{enumerate}

\item 利用对称性简化计算：设 $f(x)=x^5+x^3+x$，则 $f(2)+f(-2)=$ \score{5}
\begin{enumerate}[label=\Alph*.]
    \item $0$
    \item $40$
    \item $48$
    \item $80$
\end{enumerate}

\item 设等比数列 $\{a_n\}$ 满足 $a_1=1$，$a_n>0$，$S_n$ 为其前 $n$ 项和，$S_3=7$，则公比 $q=$ \score{5}
\begin{enumerate}[label=\Alph*.]
    \item $1$
    \item $2$
    \item $3$
    \item $-3$
\end{enumerate}

\item 设 $F(x)=\displaystyle\int_0^x f(t)\,dt$，其中 $f(t)$ 为连续函数。若 $f(x)$ 为奇函数，则 $F(x)$ 为 \score{5}
\begin{enumerate}[label=\Alph*.]
    \item 奇函数
    \item 偶函数
    \item 非奇非偶函数
    \item 视 $f$ 的具体形式而定
\end{enumerate}

\end{enumerate}

\vspace{10pt}

{\large\bfseries 二、多项选择题：本题共4小题，每小题5分，共20分。每小题有多个选项符合题目要求，全部选对得5分，部分选对得2分，有选错的得0分。}

\vspace{8pt}

\begin{enumerate}[resume]

\item 关于数形结合思想，下列说法正确的是 \score{5}
\begin{enumerate}[label=\Alph*.]
    \item 方程 $\sin x=\dfrac{x}{10}$ 的实根个数可以通过画图确定
    \item 不等式 $|x-1|+|x+2|\geq 3$ 的解集可以利用数轴上的距离解释
    \item 函数 $y=f(x)$ 与 $y=g(x)$ 的交点横坐标即为方程 $f(x)=g(x)$ 的根
    \item 数形结合只能用于解方程，不能用于证明不等式
\end{enumerate}

\item 关于分类讨论思想，下列需要分类讨论的是 \score{5}
\begin{enumerate}[label=\Alph*.]
    \item 解含参数的不等式 $ax^2+bx+c>0$
    \item 求 $|x-1|+|x-2|$ 的最小值
    \item 求等比数列前 $n$ 项和 $S_n$（需要讨论公比 $q$ 是否为 $1$）
    \item 求导数 $f'(x)$ 并确定单调区间时，需要讨论导数的正负
\end{enumerate}

\item 关于函数与方程思想，下列正确的是 \score{5}
\begin{enumerate}[label=\Alph*.]
    \item 求方程 $f(x)=0$ 的根等价于求函数 $y=f(x)$ 的零点
    \item 利用导数研究函数性质可以判断方程根的个数
    \item 构造函数是证明不等式的常用方法
    \item 函数思想的核心是建立变量之间的对应关系
\end{enumerate}

\item 关于转化与化归思想，下列正确的是 \score{5}
\begin{enumerate}[label=\Alph*.]
    \item 复杂问题可以转化为简单问题
    \item 未知问题可以转化为已知问题
    \item 立体几何问题可以通过向量方法转化为代数问题
    \item 转化过程中不需要保持等价性
\end{enumerate}

\end{enumerate}

\vspace{10pt}

{\large\bfseries 三、填空题：本题共4小题，每小题5分，共20分。}

\vspace{8pt}

\begin{enumerate}[resume]

\item 设 $f(x)=\dfrac{1}{1+2^x}$，则 $f(x)+f(-x)=$ \underline{\hspace{2cm}}。\score{5}

\item 设 $a_n=2n-1$，则 $\displaystyle\sum_{k=1}^{n}a_k=$ \underline{\hspace{2cm}}，$\displaystyle\sum_{k=1}^{n}(-1)^ka_k=$ \underline{\hspace{2cm}}。\score{5}

\item 设 $f(x)=x^3-3ax^2+3x+1$ 在 $x=1$ 处取极值，则 $a=$ \underline{\hspace{2cm}}，该极值为 \underline{\hspace{2cm}}（极大值或极小值）。\score{5}

\item 从 $1,2,\cdots,9$ 中选取 $3$ 个不同的数使其和为偶数，共有 \underline{\hspace{2cm}} 种选法。\score{5}

\end{enumerate}

\vspace{10pt}

{\large\bfseries 四、解答题：本题共6小题，共70分。解答应写出文字说明、证明过程或演算步骤。}

\vspace{8pt}

\begin{enumerate}[resume]

\item （10分）构造函数法证明不等式。

\begin{enumerate}[label=(\arabic*)]
    \item 设 $f(x)=e^x-x-1$，证明 $f(x)\geq 0$ 对所有 $x\in\mathbb{R}$ 成立（即 $e^x\geq x+1$）；
    \item 利用上述结论证明：$\ln x\leq x-1$ 对所有 $x>0$ 成立；
    \item 利用上述结论证明：对 $a,b>0$，有 $\dfrac{a+b}{2}\geq \sqrt{ab}$（提示：令 $x=\dfrac{a}{b}$）。
\end{enumerate}

\item （12分）数学归纳法的深化应用。

\begin{enumerate}[label=(\arabic*)]
    \item 用数学归纳法证明：$1^2+2^2+\cdots+n^2=\dfrac{n(n+1)(2n+1)}{6}$；
    \item 用数学归纳法证明：$1+\dfrac{1}{2^2}+\dfrac{1}{3^2}+\cdots+\dfrac{1}{n^2}<2-\dfrac{1}{n}$（$n\geq 2$）；
    \item 用数学归纳法证明：$n!>\left(\dfrac{n}{e}\right)^n$ 对所有 $n\geq 1$ 成立（提示：利用 $e^x\geq x+1$ 或 $(1+\frac{1}{n})^n<e$）。
\end{enumerate}

\item （12分）整体思想与换元法的综合应用。

\begin{enumerate}[label=(\arabic*)]
    \item 求函数 $f(x)=\dfrac{x^2+3}{\sqrt{x^2+2}}$ 的最小值（提示：设 $t=\sqrt{x^2+2}\geq \sqrt{2}$）；
    \item 设 $x+y=1$，$x,y>0$，求 $\dfrac{1}{x}+\dfrac{2}{y}$ 的最小值（整体代换法）；
    \item 求 $\displaystyle\int_0^1 x\sqrt{1-x^2}\,dx$（提示：换元 $x=\sin\theta$）。
\end{enumerate}

\item （12分）化归思想在数列问题中的应用。

\begin{enumerate}[label=(\arabic*)]
    \item 设数列 $\{a_n\}$ 满足 $a_1=1$，$a_{n+1}=2a_n+1$（$n\geq 1$）。设 $b_n=a_n+1$，证明 $\{b_n\}$ 为等比数列，并求 $a_n$；
    \item 设数列 $\{a_n\}$ 满足 $a_1=1$，$a_{n+1}=\dfrac{a_n}{2a_n+1}$。设 $b_n=\dfrac{1}{a_n}$，求 $b_n$ 并由此求 $a_n$；
    \item 设数列 $\{a_n\}$ 满足 $a_1=1$，$S_n=\dfrac{n+2}{3}a_n$，求 $a_n$（提示：用 $S_n-S_{n-1}=a_n$ 建立递推关系，再化归为 $\dfrac{a_{n+1}}{a_n}$ 的关系）。
\end{enumerate}

\item （12分）极限与逼近思想。

\begin{enumerate}[label=(\arabic*)]
    \item 设 $a_1=\sqrt{2}$，$a_{n+1}=\sqrt{2+a_n}$（$n\geq 1$）。证明 $\{a_n\}$ 单调递增且有上界 $2$，并求 $\displaystyle\lim_{n\to\infty}a_n$；
    \item 利用极限的夹逼定理求 $\displaystyle\lim_{n\to\infty}\left(\dfrac{1}{\sqrt{n^2+1}}+\dfrac{1}{\sqrt{n^2+2}}+\cdots+\dfrac{1}{\sqrt{n^2+n}}\right)$；
    \item 证明 $\displaystyle\lim_{n\to\infty}\left(1+\dfrac{1}{n}\right)^n$ 存在（提示：证明数列单调递增且有上界 $3$，利用二项式定理展开与 $\displaystyle\sum_{k=0}^{n}\dfrac{1}{k!}$ 比较）。
\end{enumerate}

\item （12分）综合运用多种数学思想解高考压轴题。

\begin{enumerate}[label=(\arabic*)]
    \item 设 $f(x)=x\ln x-ax^2+(2a-1)x$（$a\in\mathbb{R}$）。若 $f(x)$ 在 $x=1$ 处取极值，求 $a$ 的值；
    \item 在（1）的条件下，求 $f(x)$ 的单调区间；
    \item 证明：当 $n\geq 2$，$n\in\mathbb{N}^*$ 时，$\ln 2+\ln 3+\cdots+\ln n>\dfrac{n-1}{2}$（提示：利用（2）的结果或构造辅助函数）。
\end{enumerate}

\end{enumerate}

\end{document}
